% Generated draft: requires mathematical and source review.
\documentclass[11pt]{article}
\usepackage[margin=25mm]{geometry}
\usepackage{fontspec}
\setmainfont{Latin Modern Roman}
\usepackage{amsmath,amssymb,mathtools,unicode-math}
\setmathfont{Latin Modern Math}
\usepackage{enumitem,needspace,xurl,hyperref}
\setlength{\parindent}{0pt}
\setlength{\parskip}{5pt}
\setlength{\emergencystretch}{3em}
\newcommand{\N}{\mathbb{N}}
\newcommand{\Z}{\mathbb{Z}}
\newcommand{\Q}{\mathbb{Q}}
\newcommand{\R}{\mathbb{R}}
\newcommand{\C}{\mathbb{C}}
\newcommand{\F}{\mathbb{F}}
\newcommand{\A}{\mathbb{A}}
\newcommand{\PP}{\mathbb{P}}
\newcommand{\E}{\mathbb{E}}
\newcommand{\eps}{\varepsilon}
\newcommand{\dd}{\,\mathrm{d}}
\begin{document}
\section*{UnsolvedMath ID: 30003216}
\textbf{Source problem code:} catalogue entry
\subsection{Definitions and mathematical statement}
Let \(\mathcal A\) denote the adaptive hybrid finite-element procedure referred to in the supplied source record, driven by its proposed two-step a posteriori flux-error estimator. Here these terms identify the source-specific method and estimator; the supplied record does not provide their mathematical definitions.

\textbf{Question.} For the class of problems and under the hypotheses intended in the cited report, does \(\mathcal A\) converge?

The supplied record does not specify the governing partial differential equation, spatial domain, boundary conditions, coefficient or data assumptions, finite-element spaces, construction of the two-step estimator, marking rule, refinement rule, or stopping criterion. It also does not specify the quantity required to converge, its limiting value, or the topology or norm of convergence. In particular, convergence of an estimator, convergence of a flux approximation, and convergence of a discrete solution are distinct possible targets; the record does not identify the intended target precisely. Consequently, a fully quantified convergence statement cannot be recovered from the supplied evidence without additional source details.

The supplied background states that reliability and efficiency of the hybrid two-step finite-element flux estimator are known, while convergence of the adaptive procedure was explicitly left open. These are supplied literature claims, not independently verified facts in this draft; the hypotheses and precise inequalities underlying reliability and efficiency are absent.
\subsection{Short English statement}
The question asks whether the adaptive hybrid finite-element method driven by the source's two-step flux-error estimator converges. The supplied record omits the equations, algorithmic hypotheses, and meaning of convergence needed for a complete mathematical formulation. This statement-only draft therefore retains the question and explicitly records those omissions for human mathematical review.
\subsection{Sources}
\begin{itemize}
\item Local UnsolvedMath registry: \texttt{data/unsolvedmath/problems.json}, permanent ID \(30003216\), source problem code \texttt{catalogue entry}, title \emph{Convergence of Adaptive Hybrid Finite Element Methods}; collection: Oberwolfach Reports Open Problems. The supplied record gives proposed year 2016 and names no individual proposer.
\item Supplied report citation: \emph{Self-Adaptive Numerical Methods for Computationally Challenging Problems} (2017), Oberwolfach Reports, pp.\ 2399--2464, DOI \texttt{10.4171/owr/2016/42}. Original open-problem link: \url{https://doi.org/10.4171/owr/2016/42}. The record describes these reports as documenting mathematical workshops and their open-problem discussions.
\item Supplied literature source labeled ``later hybrid estimator work'': \url{https://doi.org/10.1016/S0377-0427(23)00615-5}. No author, title, or specific supporting passage is supplied for this reference.
\item Literature assessment checked at \texttt{2026-08-21}: the supplied curation labels the problem open and reports that reliability and efficiency are known, that adaptive convergence was explicitly left open, and that no general proof was located. This is supplied dated curation derived from the OWR record, not an independently verified literature assessment; time-sensitive status and cited claims require independent verification.
\item The registry assigns statement status \texttt{corrected\_verified} and reports that the official report passage was checked and rewritten while preserving its mathematical content. This describes the supplied provenance claim; the report passage and literature sources have not been independently inspected for this draft.
\end{itemize}
\end{document}
